% TOP_PROBLEM: {"id": 2793, "title": "Kirby Problem 2.45", "category_id": 7, "category": {"id": 7, "name": "topology", "display_name": "Topology", "description": "Properties preserved under continuous deformations.", "slug": "topology", "order_index": 7, "created_at": "2026-07-31T15:26:25.671Z"}, "status": "open"}
% FIRST_POSTED: null
% ENTRY_KIND: statement_only
% Imported from: batch11.tex; UnsolvedMath ID 2793
\documentclass[11pt]{article}
\usepackage[margin=25mm]{geometry}
\usepackage{fontspec}
\setmainfont{Latin Modern Roman}
\setsansfont{Latin Modern Sans}
\usepackage{amsmath,amssymb,mathtools}
\usepackage{unicode-math}
\setmathfont{Latin Modern Math}
\usepackage{xurl,hyperref}
\setcounter{secnumdepth}{0}
\setlength{\parindent}{0pt}
\setlength{\parskip}{5pt}
\setlength{\emergencystretch}{3em}
\newcommand{\sref}[2]{\hyperlink{source:#1:#2}{[S#2]}}
\newcommand{\cataloguescope}[1]{\paragraph{Recorded target.}#1}
\begin{document}
% UnsolvedMath ID 2793; source catalogue code KP-2.45
\setcounter{section}{2792}
\section{Kirby Problem 2.45}\label{2793}
{\small\sffamily UnsolvedMath ID 2793 · Topology}
\subsection{Definitions and mathematical statement}

\paragraph{Shared notation.}$\mathbb N=\{1,2,\ldots\}$, and $\mathbb Z,\mathbb Q,\mathbb R,\mathbb C$ denote the integers, rationals, reals and complex numbers. Any other conventions are specified locally.


An end-periodic oriented manifold has finitely many ends, each with a neighborhood modeled on one half of an infinite cyclic cover of a compact manifold. An end-periodic automorphism agrees, sufficiently far out on each end, with a covering translation in that model. Use either the smooth category, with diffeomorphisms, or the topological category, with homeomorphisms.

An oriented end-periodic cobordism from $(S_0,f_0)$ to $(S_1,f_1)$ is an oriented 3-manifold $W$ with $\partial W=S_0\sqcup(-S_1)$, equipped with an end-periodic automorphism restricting to $f_0,f_1$ on the boundary. All manifolds, end models and automorphisms are in the chosen category. The group $\Delta^e_2$ consists of cobordism classes of these surface-automorphism pairs; addition is disjoint union and inversion reverses orientation. The end-periodic restriction is imposed on the extension as well as on the boundary maps.
\paragraph{Statement recorded in the supplied source.}
Compute the oriented end-periodic cobordism group $\Delta^e_2$ of surface automorphisms, in the diffeomorphism and homeomorphism categories. Determine its group structure rather than substituting the ordinary compact surface-automorphism cobordism group. \sref{2793}{2}

\subsection{Short English statement}
Compute cobordism classes of end-periodic surface automorphisms, requiring the cobordism extension to remain end-periodic, in both smooth and topological categories.
\subsection{Sources}
\begin{description}
\item[\hypertarget{source:2793:1}{S1}] ULAM.AI and the UnsolvedMath Contributors, UnsolvedMath: A Curated Collection of Open Mathematics Problems, record 2793 (KP-2.45). CC BY 4.0. \url{https://huggingface.co/datasets/ulamai/UnsolvedMath}.
\item[\hypertarget{source:2793:2}{S2}] R. İnanç Baykur, Robion C. Kirby and Daniel Ruberman, eds., \emph{K3: A New Problem List in Low-Dimensional Topology}, Mathematical Surveys and Monographs 295, American Mathematical Society (2026), Problem 2.45 and its remarks; Chapters 1 and 2 for the stated conventions. \url{https://math.berkeley.edu/sites/default/files/surv-295-ruberman-watermarked-author-pdf.pdf}.
\end{description}
\label{end:2793}
\end{document}
